% titel: Blatt 1, A4: Protonen, Neutronen, Elektronen
% skript: Mortimer Kap. 2 (Nuklidbezeichnung, Ladung)

% Teil I: Werkzeuge, kein Lösungsweg.

Aus Nuklidschreibweise ${}^A_Z\mathrm{E}^{q}$ Zahl der Bausteine ableiten.

\begin{werkzeugkasten}[Bestandteile aus Nuklidsymbol]
Für ${}^A_Z\mathrm{E}^{q}$:
\begin{align}
  \text{Protonen } Z &= \text{Ordnungszahl},\\
  \text{Neutronen } N &= A - Z,\\
  \text{Elektronen} &= Z - q \quad (q=\text{Ionenladung; neutral: }q=0).
\end{align}
Bei Kationen ($q>0$) fehlen $q$ Elektronen; die Kernzusammensetzung ($Z,N$)
bleibt unverändert.

\textit{Skript:} Mortimer Kap. 2.4 (Nuklide, Isotope, Ionen), Übung 1/4.
\end{werkzeugkasten}

\begin{achtung}[Darauf achten]
\begin{itemize}
  \item Die Ordnungszahl $Z$ ist durch das Element festgelegt
        ($\mathrm{Cu}=29$, $\mathrm{Ag}=47$, $\mathrm{U}=92$, \dots{} ) --- aus dem
        Periodensystem nehmen.
  \item Bei $N=Z$ kleiner Kern, mit wachsender Masse nimmt das Neutronen-/
        Protonen-Verhältnis zu (stabilisierender Effekt gegen Coulomb-Abstossung).
  \item Kation: $2+$ bedeutet zwei Elektronen weniger, nicht mehr.
\end{itemize}
\end{achtung}
